% Options for packages loaded elsewhere
% Options for packages loaded elsewhere
\PassOptionsToPackage{unicode}{hyperref}
\PassOptionsToPackage{hyphens}{url}
\PassOptionsToPackage{dvipsnames,svgnames,x11names}{xcolor}
%
\documentclass[
  letterpaper,
  DIV=11,
  numbers=noendperiod]{scrartcl}
\usepackage{xcolor}
\usepackage[margin=1in]{geometry}
\usepackage{amsmath,amssymb}
\setcounter{secnumdepth}{-\maxdimen} % remove section numbering
\usepackage{iftex}
\ifPDFTeX
  \usepackage[T1]{fontenc}
  \usepackage[utf8]{inputenc}
  \usepackage{textcomp} % provide euro and other symbols
\else % if luatex or xetex
  \usepackage{unicode-math} % this also loads fontspec
  \defaultfontfeatures{Scale=MatchLowercase}
  \defaultfontfeatures[\rmfamily]{Ligatures=TeX,Scale=1}
\fi
\usepackage{lmodern}
\ifPDFTeX\else
  % xetex/luatex font selection
\fi
% Use upquote if available, for straight quotes in verbatim environments
\IfFileExists{upquote.sty}{\usepackage{upquote}}{}
\IfFileExists{microtype.sty}{% use microtype if available
  \usepackage[]{microtype}
  \UseMicrotypeSet[protrusion]{basicmath} % disable protrusion for tt fonts
}{}
\makeatletter
\@ifundefined{KOMAClassName}{% if non-KOMA class
  \IfFileExists{parskip.sty}{%
    \usepackage{parskip}
  }{% else
    \setlength{\parindent}{0pt}
    \setlength{\parskip}{6pt plus 2pt minus 1pt}}
}{% if KOMA class
  \KOMAoptions{parskip=half}}
\makeatother
% Make \paragraph and \subparagraph free-standing
\makeatletter
\ifx\paragraph\undefined\else
  \let\oldparagraph\paragraph
  \renewcommand{\paragraph}{
    \@ifstar
      \xxxParagraphStar
      \xxxParagraphNoStar
  }
  \newcommand{\xxxParagraphStar}[1]{\oldparagraph*{#1}\mbox{}}
  \newcommand{\xxxParagraphNoStar}[1]{\oldparagraph{#1}\mbox{}}
\fi
\ifx\subparagraph\undefined\else
  \let\oldsubparagraph\subparagraph
  \renewcommand{\subparagraph}{
    \@ifstar
      \xxxSubParagraphStar
      \xxxSubParagraphNoStar
  }
  \newcommand{\xxxSubParagraphStar}[1]{\oldsubparagraph*{#1}\mbox{}}
  \newcommand{\xxxSubParagraphNoStar}[1]{\oldsubparagraph{#1}\mbox{}}
\fi
\makeatother

\usepackage{color}
\usepackage{fancyvrb}
\newcommand{\VerbBar}{|}
\newcommand{\VERB}{\Verb[commandchars=\\\{\}]}
\DefineVerbatimEnvironment{Highlighting}{Verbatim}{commandchars=\\\{\}}
% Add ',fontsize=\small' for more characters per line
\usepackage{framed}
\definecolor{shadecolor}{RGB}{241,243,245}
\newenvironment{Shaded}{\begin{snugshade}}{\end{snugshade}}
\newcommand{\AlertTok}[1]{\textcolor[rgb]{0.68,0.00,0.00}{#1}}
\newcommand{\AnnotationTok}[1]{\textcolor[rgb]{0.37,0.37,0.37}{#1}}
\newcommand{\AttributeTok}[1]{\textcolor[rgb]{0.40,0.45,0.13}{#1}}
\newcommand{\BaseNTok}[1]{\textcolor[rgb]{0.68,0.00,0.00}{#1}}
\newcommand{\BuiltInTok}[1]{\textcolor[rgb]{0.00,0.23,0.31}{#1}}
\newcommand{\CharTok}[1]{\textcolor[rgb]{0.13,0.47,0.30}{#1}}
\newcommand{\CommentTok}[1]{\textcolor[rgb]{0.37,0.37,0.37}{#1}}
\newcommand{\CommentVarTok}[1]{\textcolor[rgb]{0.37,0.37,0.37}{\textit{#1}}}
\newcommand{\ConstantTok}[1]{\textcolor[rgb]{0.56,0.35,0.01}{#1}}
\newcommand{\ControlFlowTok}[1]{\textcolor[rgb]{0.00,0.23,0.31}{\textbf{#1}}}
\newcommand{\DataTypeTok}[1]{\textcolor[rgb]{0.68,0.00,0.00}{#1}}
\newcommand{\DecValTok}[1]{\textcolor[rgb]{0.68,0.00,0.00}{#1}}
\newcommand{\DocumentationTok}[1]{\textcolor[rgb]{0.37,0.37,0.37}{\textit{#1}}}
\newcommand{\ErrorTok}[1]{\textcolor[rgb]{0.68,0.00,0.00}{#1}}
\newcommand{\ExtensionTok}[1]{\textcolor[rgb]{0.00,0.23,0.31}{#1}}
\newcommand{\FloatTok}[1]{\textcolor[rgb]{0.68,0.00,0.00}{#1}}
\newcommand{\FunctionTok}[1]{\textcolor[rgb]{0.28,0.35,0.67}{#1}}
\newcommand{\ImportTok}[1]{\textcolor[rgb]{0.00,0.46,0.62}{#1}}
\newcommand{\InformationTok}[1]{\textcolor[rgb]{0.37,0.37,0.37}{#1}}
\newcommand{\KeywordTok}[1]{\textcolor[rgb]{0.00,0.23,0.31}{\textbf{#1}}}
\newcommand{\NormalTok}[1]{\textcolor[rgb]{0.00,0.23,0.31}{#1}}
\newcommand{\OperatorTok}[1]{\textcolor[rgb]{0.37,0.37,0.37}{#1}}
\newcommand{\OtherTok}[1]{\textcolor[rgb]{0.00,0.23,0.31}{#1}}
\newcommand{\PreprocessorTok}[1]{\textcolor[rgb]{0.68,0.00,0.00}{#1}}
\newcommand{\RegionMarkerTok}[1]{\textcolor[rgb]{0.00,0.23,0.31}{#1}}
\newcommand{\SpecialCharTok}[1]{\textcolor[rgb]{0.37,0.37,0.37}{#1}}
\newcommand{\SpecialStringTok}[1]{\textcolor[rgb]{0.13,0.47,0.30}{#1}}
\newcommand{\StringTok}[1]{\textcolor[rgb]{0.13,0.47,0.30}{#1}}
\newcommand{\VariableTok}[1]{\textcolor[rgb]{0.07,0.07,0.07}{#1}}
\newcommand{\VerbatimStringTok}[1]{\textcolor[rgb]{0.13,0.47,0.30}{#1}}
\newcommand{\WarningTok}[1]{\textcolor[rgb]{0.37,0.37,0.37}{\textit{#1}}}

\usepackage{longtable,booktabs,array}
\newcounter{none} % for unnumbered tables
\usepackage{calc} % for calculating minipage widths
% Correct order of tables after \paragraph or \subparagraph
\usepackage{etoolbox}
\makeatletter
\patchcmd\longtable{\par}{\if@noskipsec\mbox{}\fi\par}{}{}
\makeatother
% Allow footnotes in longtable head/foot
\IfFileExists{footnotehyper.sty}{\usepackage{footnotehyper}}{\usepackage{footnote}}
\makesavenoteenv{longtable}
\usepackage{graphicx}
\makeatletter
\newsavebox\pandoc@box
\newcommand*\pandocbounded[1]{% scales image to fit in text height/width
  \sbox\pandoc@box{#1}%
  \Gscale@div\@tempa{\textheight}{\dimexpr\ht\pandoc@box+\dp\pandoc@box\relax}%
  \Gscale@div\@tempb{\linewidth}{\wd\pandoc@box}%
  \ifdim\@tempb\p@<\@tempa\p@\let\@tempa\@tempb\fi% select the smaller of both
  \ifdim\@tempa\p@<\p@\scalebox{\@tempa}{\usebox\pandoc@box}%
  \else\usebox{\pandoc@box}%
  \fi%
}
% Set default figure placement to htbp
\def\fps@figure{htbp}
\makeatother





\setlength{\emergencystretch}{3em} % prevent overfull lines

\providecommand{\tightlist}{%
  \setlength{\itemsep}{0pt}\setlength{\parskip}{0pt}}



 


\KOMAoption{captions}{tableheading}
\makeatletter
\@ifpackageloaded{tcolorbox}{}{\usepackage[skins,breakable]{tcolorbox}}
\@ifpackageloaded{fontawesome5}{}{\usepackage{fontawesome5}}
\definecolor{quarto-callout-color}{HTML}{909090}
\definecolor{quarto-callout-note-color}{HTML}{0758E5}
\definecolor{quarto-callout-important-color}{HTML}{CC1914}
\definecolor{quarto-callout-warning-color}{HTML}{EB9113}
\definecolor{quarto-callout-tip-color}{HTML}{00A047}
\definecolor{quarto-callout-caution-color}{HTML}{FC5300}
\definecolor{quarto-callout-color-frame}{HTML}{acacac}
\definecolor{quarto-callout-note-color-frame}{HTML}{4582ec}
\definecolor{quarto-callout-important-color-frame}{HTML}{d9534f}
\definecolor{quarto-callout-warning-color-frame}{HTML}{f0ad4e}
\definecolor{quarto-callout-tip-color-frame}{HTML}{02b875}
\definecolor{quarto-callout-caution-color-frame}{HTML}{fd7e14}
\makeatother
\makeatletter
\@ifpackageloaded{caption}{}{\usepackage{caption}}
\AtBeginDocument{%
\ifdefined\contentsname
  \renewcommand*\contentsname{Table of contents}
\else
  \newcommand\contentsname{Table of contents}
\fi
\ifdefined\listfigurename
  \renewcommand*\listfigurename{List of Figures}
\else
  \newcommand\listfigurename{List of Figures}
\fi
\ifdefined\listtablename
  \renewcommand*\listtablename{List of Tables}
\else
  \newcommand\listtablename{List of Tables}
\fi
\ifdefined\figurename
  \renewcommand*\figurename{Figure}
\else
  \newcommand\figurename{Figure}
\fi
\ifdefined\tablename
  \renewcommand*\tablename{Table}
\else
  \newcommand\tablename{Table}
\fi
}
\@ifpackageloaded{float}{}{\usepackage{float}}
\floatstyle{ruled}
\@ifundefined{c@chapter}{\newfloat{codelisting}{h}{lop}}{\newfloat{codelisting}{h}{lop}[chapter]}
\floatname{codelisting}{Listing}
\newcommand*\listoflistings{\listof{codelisting}{List of Listings}}
\makeatother
\makeatletter
\makeatother
\makeatletter
\@ifpackageloaded{caption}{}{\usepackage{caption}}
\@ifpackageloaded{subcaption}{}{\usepackage{subcaption}}
\makeatother
\usepackage{bookmark}
\IfFileExists{xurl.sty}{\usepackage{xurl}}{} % add URL line breaks if available
\urlstyle{same}
% fallback for those not using the hyperref driver hyperxmp:
\makeatletter
\@ifundefined{xmpquote}{\newcommand{\xmpquote}[1]{#1}}{}
\makeatother
\hypersetup{
  pdftitle={W06 Practice: Interpreting an Ordinal Path Model},
  pdfauthor={Susanne DeVore, PhD},
  colorlinks=true,
  linkcolor={blue},
  filecolor={Maroon},
  citecolor={Blue},
  urlcolor={Blue},
  pdfcreator={LaTeX via pandoc}}


\title{W06 Practice: Interpreting an Ordinal Path Model}
\usepackage{etoolbox}
\makeatletter
\providecommand{\subtitle}[1]{% add subtitle to \maketitle
  \apptocmd{\@title}{\par {\large #1 \par}}{}{}
}
\makeatother
\subtitle{EDEP 625 Structural Equation Modeling}
\author{Susanne DeVore, PhD}
\date{2026-10-01}
\begin{document}
\maketitle


\subsection{Purpose}\label{purpose}

Practice reporting on an ordinal MVPA.

\begin{tcolorbox}[enhanced jigsaw, arc=.35mm, bottomrule=.15mm, breakable, colback=white, colframe=quarto-callout-note-color-frame, left=2mm, leftrule=.75mm, opacityback=0, rightrule=.15mm, toprule=.15mm]
\begin{minipage}[t]{5.5mm}
\textcolor{quarto-callout-note-color}{\faInfo}
\end{minipage}%
\begin{minipage}[t]{\textwidth - 5.5mm}

\textbf{This is all fake data}, generated specifically to illustrate an
ordinal measured-variable path analysis.

\end{minipage}%
\end{tcolorbox}

\subsection{The scenario}\label{the-scenario}

\textbf{Research question.} Among undergraduates in online courses, do
\textbf{instructor responsiveness} and \textbf{peer interaction} relate
to students' \textbf{intention to persist} in the course \emph{through}
their \textbf{sense of belonging}?

\textbf{The data.} A survey of 450 students. Every variable is a Likert
item.

{\def\LTcaptype{none} % do not increment counter
\begin{longtable}[]{@{}
  >{\raggedright\arraybackslash}p{(\linewidth - 4\tabcolsep) * \real{0.1183}}
  >{\raggedright\arraybackslash}p{(\linewidth - 4\tabcolsep) * \real{0.5269}}
  >{\raggedright\arraybackslash}p{(\linewidth - 4\tabcolsep) * \real{0.3548}}@{}}
\toprule\noalign{}
\begin{minipage}[b]{\linewidth}\raggedright
Variable
\end{minipage} & \begin{minipage}[b]{\linewidth}\raggedright
Meaning
\end{minipage} & \begin{minipage}[b]{\linewidth}\raggedright
Scale
\end{minipage} \\
\midrule\noalign{}
\endhead
\bottomrule\noalign{}
\endlastfoot
\texttt{respond} & ``My instructor responds to questions promptly'' &
1--4 (strongly disagree → agree) \\
\texttt{peer} & ``I interact regularly with classmates'' & 1--4
(strongly disagree → agree) \\
\texttt{belong} & sense of belonging in the course & 1--5 (not at all →
completely) \\
\texttt{persist} & intention to continue / re-enroll & 1--4 (very
unlikely → very likely) \\
\end{longtable}
}

\textbf{The hypothesized model.}

\begin{center}
\pandocbounded{\includegraphics[keepaspectratio]{w06_mvpa_ordinal_practice_files/figure-pdf/unnamed-chunk-1-1.pdf}}
\end{center}

\subsection{The data}\label{the-data}

\begin{Shaded}
\begin{Highlighting}[]
\FunctionTok{library}\NormalTok{(lavaan)}

\FunctionTok{set.seed}\NormalTok{(}\DecValTok{625}\NormalTok{)}
\NormalTok{n }\OtherTok{\textless{}{-}} \DecValTok{450}
\NormalTok{respond\_s }\OtherTok{\textless{}{-}} \FunctionTok{rnorm}\NormalTok{(n)}
\NormalTok{peer\_s    }\OtherTok{\textless{}{-}} \FloatTok{0.30} \SpecialCharTok{*}\NormalTok{ respond\_s }\SpecialCharTok{+} \FunctionTok{rnorm}\NormalTok{(n, }\DecValTok{0}\NormalTok{, }\FunctionTok{sqrt}\NormalTok{(}\DecValTok{1} \SpecialCharTok{{-}} \FloatTok{0.30}\SpecialCharTok{\^{}}\DecValTok{2}\NormalTok{))}
\NormalTok{belong\_s  }\OtherTok{\textless{}{-}} \FloatTok{0.45} \SpecialCharTok{*}\NormalTok{ respond\_s }\SpecialCharTok{+} \FloatTok{0.35} \SpecialCharTok{*}\NormalTok{ peer\_s }\SpecialCharTok{+}
             \FunctionTok{rnorm}\NormalTok{(n, }\DecValTok{0}\NormalTok{, }\FunctionTok{sqrt}\NormalTok{(}\DecValTok{1} \SpecialCharTok{{-}} \FloatTok{0.45}\SpecialCharTok{\^{}}\DecValTok{2} \SpecialCharTok{{-}} \FloatTok{0.35}\SpecialCharTok{\^{}}\DecValTok{2}\NormalTok{))}
\NormalTok{persist\_s }\OtherTok{\textless{}{-}} \FloatTok{0.50} \SpecialCharTok{*}\NormalTok{ belong\_s }\SpecialCharTok{+} \FunctionTok{rnorm}\NormalTok{(n, }\DecValTok{0}\NormalTok{, }\FunctionTok{sqrt}\NormalTok{(}\DecValTok{1} \SpecialCharTok{{-}} \FloatTok{0.50}\SpecialCharTok{\^{}}\DecValTok{2}\NormalTok{))}

\NormalTok{to4 }\OtherTok{\textless{}{-}} \ControlFlowTok{function}\NormalTok{(z) }\FunctionTok{as.integer}\NormalTok{(}\FunctionTok{cut}\NormalTok{(z, }\FunctionTok{c}\NormalTok{(}\SpecialCharTok{{-}}\ConstantTok{Inf}\NormalTok{, }\SpecialCharTok{{-}}\FloatTok{0.8}\NormalTok{, }\DecValTok{0}\NormalTok{, }\FloatTok{0.8}\NormalTok{, }\ConstantTok{Inf}\NormalTok{), }\AttributeTok{labels =} \DecValTok{1}\SpecialCharTok{:}\DecValTok{4}\NormalTok{))}
\NormalTok{to5 }\OtherTok{\textless{}{-}} \ControlFlowTok{function}\NormalTok{(z) }\FunctionTok{as.integer}\NormalTok{(}\FunctionTok{cut}\NormalTok{(z, }\FunctionTok{c}\NormalTok{(}\SpecialCharTok{{-}}\ConstantTok{Inf}\NormalTok{, }\SpecialCharTok{{-}}\DecValTok{1}\NormalTok{, }\SpecialCharTok{{-}}\FloatTok{0.2}\NormalTok{, }\FloatTok{0.5}\NormalTok{, }\FloatTok{1.1}\NormalTok{, }\ConstantTok{Inf}\NormalTok{), }\AttributeTok{labels =} \DecValTok{1}\SpecialCharTok{:}\DecValTok{5}\NormalTok{))}

\NormalTok{belong\_survey }\OtherTok{\textless{}{-}} \FunctionTok{data.frame}\NormalTok{(}
  \AttributeTok{respond =} \FunctionTok{to4}\NormalTok{(respond\_s),}
  \AttributeTok{peer    =} \FunctionTok{to4}\NormalTok{(peer\_s),}
  \AttributeTok{belong  =} \FunctionTok{to5}\NormalTok{(belong\_s),}
  \AttributeTok{persist =} \FunctionTok{to4}\NormalTok{(persist\_s)}
\NormalTok{)}

\FunctionTok{head}\NormalTok{(belong\_survey)}
\end{Highlighting}
\end{Shaded}

\begin{verbatim}
  respond peer belong persist
1       2    1      4       2
2       3    4      5       3
3       2    1      2       3
4       2    1      1       3
5       3    3      2       3
6       1    2      2       2
\end{verbatim}

\begin{Shaded}
\begin{Highlighting}[]
\FunctionTok{sapply}\NormalTok{(belong\_survey, table)   }\CommentTok{\# a few discrete categories each}
\end{Highlighting}
\end{Shaded}

\begin{verbatim}
$respond

  1   2   3   4 
 92 144 130  84 

$peer

  1   2   3   4 
 74 128 151  97 

$belong

  1   2   3   4   5 
 62 122 110  88  68 

$persist

  1   2   3   4 
 89 142 136  83 
\end{verbatim}

\subsection{The fitted model}\label{the-fitted-model}

Notice some additional arguments to last week's Rlab:

\begin{Shaded}
\begin{Highlighting}[]
\NormalTok{model }\OtherTok{\textless{}{-}} \StringTok{"}
\StringTok{  belong  \textasciitilde{} a1*respond + a2*peer}
\StringTok{  persist \textasciitilde{} b*belong}

\StringTok{  ind\_respond := a1*b}
\StringTok{  ind\_peer    := a2*b}
\StringTok{"}

\NormalTok{fit }\OtherTok{\textless{}{-}} \FunctionTok{lavaan}\NormalTok{(model, }\AttributeTok{data =}\NormalTok{ belong\_survey,}
              \AttributeTok{ordered =} \FunctionTok{c}\NormalTok{(}\StringTok{"respond"}\NormalTok{, }\StringTok{"peer"}\NormalTok{, }\StringTok{"belong"}\NormalTok{, }\StringTok{"persist"}\NormalTok{),}
              \AttributeTok{estimator =} \StringTok{"DWLS"}\NormalTok{,}
              \AttributeTok{auto.var =} \ConstantTok{TRUE}\NormalTok{,   }\CommentTok{\# disturbance variances of the endogenous variables}
              \AttributeTok{auto.cov.y =} \ConstantTok{FALSE}\NormalTok{,}\CommentTok{\# keep endogenous disturbances uncorrelated (recursive)}
              \AttributeTok{auto.th =} \ConstantTok{TRUE}\NormalTok{,    }\CommentTok{\# thresholds of the ordered variables as free parameters}
              \AttributeTok{auto.delta =} \ConstantTok{TRUE}\NormalTok{, }\CommentTok{\# scaling parameters (delta) for ordered variables}
              \AttributeTok{se =} \StringTok{"bootstrap"}\NormalTok{, }\AttributeTok{bootstrap =} \DecValTok{200}
\NormalTok{              )}
\end{Highlighting}
\end{Shaded}

\subsection{The output (path
coefficients)}\label{the-output-path-coefficients}

\begin{Shaded}
\begin{Highlighting}[]
\FunctionTok{summary}\NormalTok{(fit, }\AttributeTok{standardized =} \ConstantTok{TRUE}\NormalTok{, }\AttributeTok{fit.measures =} \ConstantTok{TRUE}\NormalTok{, }\AttributeTok{rsquare =} \ConstantTok{TRUE}\NormalTok{)}
\end{Highlighting}
\end{Shaded}

\begin{verbatim}
lavaan 0.7-2 ended normally after 8 iterations

  Estimator                                       DWLS
  Optimization method                           NLMINB
  Number of model parameters                        10

  Number of observations                           450

Model Test User Model:
                                                      
  Test statistic                                 2.224
  Degrees of freedom                                 2
  P-value (Unknown)                                 NA

Model Test Baseline Model:

  Test statistic                                95.633
  Degrees of freedom                                 1
  P-value                                           NA

User Model versus Baseline Model:

  Comparative Fit Index (CFI)                    0.998
  Tucker-Lewis Index (TLI)                       0.999

Root Mean Square Error of Approximation:

  RMSEA                                          0.016
  90 Percent confidence interval - lower         0.000
  90 Percent confidence interval - upper         0.096
  P-value H_0: RMSEA <= 0.050                    0.647
  P-value H_0: RMSEA >= 0.080                    0.118

Standardized Root Mean Square Residual:

  SRMR                                           0.013

Goodness of Fit Index:

  Goodness of Fit Index (GFI)                    1.000
  90 Percent confidence interval - lower         0.982
  90 Percent confidence interval - upper         1.000

Parameter Estimates:

  Parameterization                               Delta
  Standard errors                            Bootstrap
  Number of requested bootstrap draws              200
  Number of successful bootstrap draws             200

Regressions:
                   Estimate  Std.Err  z-value  P(>|z|)   Std.lv  Std.all
  belong ~                                                              
    respond   (a1)    0.476    0.060    7.936    0.000    0.476    0.408
    peer      (a2)    0.326    0.055    5.894    0.000    0.326    0.275
  persist ~                                                             
    belong     (b)    0.407    0.045    9.048    0.000    0.407    0.468

Thresholds:
                   Estimate  Std.Err  z-value  P(>|z|)   Std.lv  Std.all
    belong|t1         0.803    0.178    4.520    0.000    0.803    0.676
    belong|t2         1.837    0.189    9.728    0.000    1.837    1.546
    belong|t3         2.598    0.208   12.465    0.000    2.598    2.187
    belong|t4         3.374    0.223   15.138    0.000    3.374    2.840
    persist|t1       -0.262    0.182   -1.439    0.150   -0.262   -0.254
    persist|t2        0.637    0.185    3.438    0.001    0.637    0.616
    persist|t3        1.521    0.194    7.850    0.000    1.521    1.472

Variances:
                   Estimate  Std.Err  z-value  P(>|z|)   Std.lv  Std.all
   .belong            1.000                               1.000    0.709
   .persist           0.835                               0.835    0.781

R-Square:
                   Estimate
    belong            0.291
    persist           0.219

Defined Parameters:
                   Estimate  Std.Err  z-value  P(>|z|)   Std.lv  Std.all
    ind_respond       0.194    0.029    6.650    0.000    0.194    0.191
    ind_peer          0.133    0.026    5.089    0.000    0.133    0.128
\end{verbatim}

\begin{Shaded}
\begin{Highlighting}[]
\CommentTok{\# Standardized estimates with standard errors, z, p, and CIs}
\FunctionTok{standardizedSolution}\NormalTok{(fit)}
\end{Highlighting}
\end{Shaded}

\begin{verbatim}
           lhs  op     rhs       label est.std    se      z pvalue ci.lower
1       belong   ~ respond          a1   0.408 0.043  9.501  0.000    0.318
2       belong   ~    peer          a2   0.275 0.043  6.333  0.000    0.186
3      persist   ~  belong           b   0.468 0.047  9.984  0.000    0.361
4       belong   |      t1               0.676 0.134  5.037  0.000    0.361
5       belong   |      t2               1.546 0.123 12.601  0.000    1.267
6       belong   |      t3               2.187 0.123 17.715  0.000    1.927
7       belong   |      t4               2.840 0.124 22.916  0.000    2.569
8      persist   |      t1              -0.254 0.177 -1.432  0.152   -0.613
9      persist   |      t2               0.616 0.176  3.506  0.000    0.230
10     persist   |      t3               1.472 0.181  8.153  0.000    1.089
11      belong  ~~  belong               0.709 0.038 18.497  0.000    0.633
12     persist  ~~ persist               0.781 0.044 17.723  0.000    0.682
13     respond  ~~ respond               1.000 0.000     NA     NA    1.000
14     respond  ~~    peer               0.223 0.000     NA     NA    0.223
15        peer  ~~    peer               1.000 0.000     NA     NA    1.000
16      belong ~*~  belong               1.000 0.000     NA     NA    1.000
17     persist ~*~ persist               1.000 0.000     NA     NA    1.000
18      belong  ~1                       0.000 0.000     NA     NA    0.000
19     persist  ~1                       0.000 0.000     NA     NA    0.000
20     respond  ~1                       2.419 0.000     NA     NA    2.419
21        peer  ~1                       2.600 0.000     NA     NA    2.600
22 ind_respond  :=    a1*b ind_respond   0.191 0.027  7.002  0.000    0.140
23    ind_peer  :=    a2*b    ind_peer   0.128 0.025  5.233  0.000    0.085
   ci.upper
1     0.484
2     0.359
3     0.564
4     0.939
5     1.776
6     2.422
7     3.100
8     0.129
9     0.996
10    1.845
11    0.787
12    0.870
13    1.000
14    0.223
15    1.000
16    1.000
17    1.000
18    0.000
19    0.000
20    2.419
21    2.600
22    0.246
23    0.183
\end{verbatim}

\begin{Shaded}
\begin{Highlighting}[]
\CommentTok{\# Unstandardized estimates with percentile bootstrap CIs}
\FunctionTok{parameterEstimates}\NormalTok{(fit, }\AttributeTok{boot.ci.type =} \StringTok{"perc"}\NormalTok{, }\AttributeTok{standardized =} \ConstantTok{TRUE}\NormalTok{)}
\end{Highlighting}
\end{Shaded}

\begin{verbatim}
           lhs  op     rhs       label    est    se      z pvalue ci.lower
1       belong   ~ respond          a1  0.476 0.060  7.936  0.000    0.355
2       belong   ~    peer          a2  0.326 0.055  5.894  0.000    0.211
3      persist   ~  belong           b  0.407 0.045  9.048  0.000    0.306
4       belong   |      t1              0.803 0.178  4.520  0.000    0.411
5       belong   |      t2              1.837 0.189  9.728  0.000    1.440
6       belong   |      t3              2.598 0.208 12.465  0.000    2.176
7       belong   |      t4              3.374 0.223 15.138  0.000    2.939
8      persist   |      t1             -0.262 0.182 -1.439  0.150   -0.626
9      persist   |      t2              0.637 0.185  3.438  0.001    0.236
10     persist   |      t3              1.521 0.194  7.850  0.000    1.108
11      belong  ~~  belong              1.000 0.000     NA     NA    1.000
12     persist  ~~ persist              0.835 0.000     NA     NA    0.835
13     respond  ~~ respond              1.033 0.000     NA     NA    1.033
14     respond  ~~    peer              0.227 0.000     NA     NA    0.227
15        peer  ~~    peer              1.002 0.000     NA     NA    1.002
16      belong ~*~  belong              1.000 0.000     NA     NA    1.000
17     persist ~*~ persist              1.000 0.000     NA     NA    1.000
18      belong  ~1                      0.000 0.000     NA     NA    0.000
19     persist  ~1                      0.000 0.000     NA     NA    0.000
20     respond  ~1                      2.458 0.000     NA     NA    2.458
21        peer  ~1                      2.602 0.000     NA     NA    2.602
22 ind_respond  :=    a1*b ind_respond  0.194 0.029  6.650  0.000    0.142
23    ind_peer  :=    a2*b    ind_peer  0.133 0.026  5.089  0.000    0.087
   ci.upper std.lv std.all std.nox
1     0.593  0.476   0.408   0.401
2     0.439  0.326   0.275   0.274
3     0.497  0.407   0.468   0.468
4     1.177  0.803   0.676   0.676
5     2.230  1.837   1.546   1.546
6     3.019  2.598   2.187   2.187
7     3.839  3.374   2.840   2.840
8     0.134 -0.262  -0.254  -0.254
9     1.037  0.637   0.616   0.616
10    1.942  1.521   1.472   1.472
11    1.000  1.000   0.709   0.709
12    0.835  0.835   0.781   0.781
13    1.033  1.033   1.000   1.033
14    0.227  0.227   0.223   0.227
15    1.002  1.002   1.000   1.002
16    1.000  1.000   1.000   1.000
17    1.000  1.000   1.000   1.000
18    0.000  0.000   0.000   0.000
19    0.000  0.000   0.000   0.000
20    2.458  2.458   2.419   2.458
21    2.602  2.602   2.600   2.602
22    0.253  0.194   0.191   0.188
23    0.191  0.133   0.128   0.128
\end{verbatim}

\subsection{Your task}\label{your-task}

\begin{enumerate}
\def\labelenumi{\arabic{enumi}.}
\tightlist
\item
  What info is needed to answer the question? (curate)
\item
  Write up a results section. Bullet points is fine. Include key
  information.
\item
  Answer the question based on the results.
\end{enumerate}




\end{document}
